\section{Association, dissociation, and rebinding}
\label{sec:binding}

In stochastic kinetics, a \emph{hazard} is the instantaneous conditional rate
of an event, given that the event has not already occurred.  If $T$ is the
waiting time to an event, its hazard is
\begin{equation}
 h(t)=\lim_{\delta t\rightarrow0}
 \frac{\Prb\!\left(t\leq T<t+\delta t\mid T\geq t\right)}{\delta t}.
 \label{eq:hazard-definition}
\end{equation}
It has units of inverse time.  For a constant hazard $h$, the probability of
no event during an interval of length $t$ is $e^{-ht}$, and the probability of
at least one event is $1-e^{-ht}$.  Thus, ``hazard'' here simply means the
instantaneous stochastic event rate; it does not imply a separate physical
mechanism.  This rate-to-probability conversion is the same one used in
chemical master-equation and kinetic Monte Carlo methods
\citep{gillespie1976}.

\subsection{Microscopic association probability}

The chemical step is
\begin{equation}
 \ce{L + R <=>[k_{on}][k_{off}] LR},
 \label{eq:reaction}
\end{equation}
where $\kOn$ is supplied in \si{\liter\per\mole\per\second} (equivalently,
$\mathrm{M}^{-1}\,\mathrm{s}^{-1}$) and $\kOff$ in
\si{\per\second}.  A ligand can react when it occupies the fluid voxel adjacent
to a receptor-bearing surface face.  If the configured reaction region
contains $n_{\mathrm{rxn}}$ voxels, its volume in liters is
\begin{equation}
 V_{\mathrm{rxn}}
 =n_{\mathrm{rxn}}a^3\left(10^3\ \mathrm{L\,m^{-3}}\right).
 \label{eq:reaction-volume}
\end{equation}
One ligand in this volume corresponds to molar concentration
\begin{equation}
 C_1=\frac{1}{\NA V_{\mathrm{rxn}}},
 \label{eq:one-particle-concentration}
\end{equation}
where $\NA$ is Avogadro's constant.  For $m$ available receptors at the same
reaction site, elementary mass action gives the total instantaneous
association rate (hazard)
\begin{equation}
 h_{\mathrm{on}}(m)=m\kOn C_1
 =\frac{m\kOn}{\NA V_{\mathrm{rxn}}}.
 \label{eq:on-hazard}
\end{equation}
Treating reaction events as a Poisson process during one update yields
\begin{equation}
 p_{\mathrm{bind}}(m)
 =1-\exp[-h_{\mathrm{on}}(m)\dt]
 =1-\exp\!\left[-\frac{m\kOn\dt}{\NA V_{\mathrm{rxn}}}\right].
 \label{eq:binding-probability}
\end{equation}
Conditional on association, one of the $m$ available receptors is selected
uniformly.  The exponential form preserves $0\le p_{\mathrm{bind}}\le1$ and
reduces to the familiar chemical-master-equation propensity
$p_{\mathrm{bind}}\simeq m\kOn\dt/(\NA V_{\mathrm{rxn}})$ when the exponent is
small \citep{gillespie1976,fange2006}.

Equation~\ref{eq:binding-probability} defines how the reported intrinsic
$\kOn$ is interpreted at the chosen lattice scale.  In a diffusion-influenced
reaction, the macroscopic association rate also depends on transport and on
the chosen reactive length scale; this distinction is familiar from
Smoluchowski and radiation-boundary descriptions
\citep{collins1949,agmon1990,andrews2009}.  Moreover, a strictly local
reaction--diffusion master equation does not possess a mesh-independent
microscopic limit for every bimolecular system
\citep{isaacson2008,hellander2012}.  Accordingly, $a$ and
$n_{\mathrm{rxn}}$ are declared model parameters, and lattice-spacing
robustness is tested explicitly rather than absorbed into an unreported rate
renormalization.

\subsection{Dissociation probability and event-wise acceleration}

For a bound complex with first-order dissociation constant $\kOff$, the
continuous waiting-time density is
\begin{equation}
 f_B(t)=\kOff e^{-\kOff t},
 \qquad \E[B]=\frac{1}{\kOff}.
 \label{eq:exponential-dwell}
\end{equation}
The dissociation hazard is the constant first-order rate $h_{\mathrm{off}}=
\kOff$.  Integrating this rate over one update gives the exact conditional
dissociation probability
\begin{equation}
 p_{\mathrm{off}}=1-e^{-\kOff\dt}.
 \label{eq:off-probability}
\end{equation}
Rather than evaluating Eq.~\ref{eq:off-probability} repeatedly during long
bound intervals, the escape assay draws the number of bound updates directly,
\begin{equation}
 G\sim\operatorname{Geometric}(p_{\mathrm{off}}),\qquad B_{\Delta}=G\dt,
 \label{eq:geometric-dwell}
\end{equation}
with support $G=1,2,\ldots$.  This event-wise skip is statistically identical
to testing dissociation after every lattice update.  Its discrete mean is
\begin{equation}
 \E[B_{\Delta}]
 =\frac{\dt}{1-e^{-\kOff\dt}}
 =\frac{1}{\kOff}+\frac{\dt}{2}+O(\kOff\dt^2),
 \label{eq:discrete-dwell-mean}
\end{equation}
so the physical exponential mean is recovered as $\kOff\dt\rightarrow0$.
This acceleration changes computational cost but not the discrete stochastic
model.

Each receptor-conditioned escape trajectory begins in the bound state.  Its
first interval is therefore drawn from Eq.~\ref{eq:geometric-dwell} and
contributes to $T_B$, but it is not counted as a rebinding event.  Upon
dissociation, the ligand is placed in the fluid voxel adjacent to the released
receptor.  It subsequently diffuses according to Section~\ref{sec:diffusion}
and can bind again according to Eq.~\ref{eq:binding-probability}.

\subsection{Self-rebinding, cross-rebinding, and finite occupancy}

A new binding event is classified relative to the receptor vacated at the
start of the current free excursion.  Return to that receptor is
\emph{self-rebinding}; binding to a different receptor is
\emph{cross-rebinding}.  The all-rebinding model permits both.  A mechanistic
self-only control disables cross-rebinding while leaving diffusion,
$\kOn$, and $\kOff$ unchanged.

Finite ensemble occupancy is represented without adding other explicitly
diffusing ligands.  For a prescribed background occupancy $\phi$, every
non-source receptor is independently available with probability $1-\phi$
during a tagged ligand's free excursion.  That availability field remains
fixed throughout the excursion and is resampled after the next dissociation;
it therefore does not flicker on the diffusion-update timescale.  The receptor
just vacated by the tagged ligand is forced available, because its ligand is
the tagged molecule itself.  This construction preserves self-return while
progressively suppressing cross-rebinding as the background approaches
saturation.

The background field is a controlled mean-field approximation, not an
explicit many-ligand equilibrium simulation.  It assumes statistically
independent receptor occupancy outside the tagged trajectory and neglects
ligand--ligand interactions and correlated turnover of background ligands.
Its advantage is that it isolates how receptor availability modifies the
single-ligand retention kernel at a known occupancy.
